%% =====================================================================
%%  fhe-handout.tex  --  A short handout on Fully Homomorphic Encryption
%%  Build: pdflatex fhe-handout (twice)
%% =====================================================================
\documentclass[11pt,a4paper]{article}

\usepackage[T1]{fontenc}
\usepackage[utf8]{inputenc}
\usepackage[sc,osf]{mathpazo}
\usepackage[scaled=0.90]{helvet}
\usepackage{microtype}
\usepackage[a4paper,top=22mm,bottom=24mm,left=24mm,right=24mm]{geometry}
\usepackage{amsmath,amssymb,mathtools,bm}
\usepackage{xcolor}
\usepackage{tikz-cd}
\usepackage{booktabs,tabularx}
\usepackage[shortlabels]{enumitem}
\setlist{itemsep=1pt,topsep=3pt}
\usepackage[most]{tcolorbox}
\usepackage{titlesec}
\usepackage{fancyhdr}
\usepackage[colorlinks,linkcolor=fheblue,citecolor=fheteal,urlcolor=fheteal]{hyperref}

%% ---- palette (same as the main text) ---------------------------------
\definecolor{fheblue}{HTML}{1F4E79}
\definecolor{fheteal}{HTML}{0F766E}
\definecolor{fheamber}{HTML}{B45309}
\definecolor{fhered}{HTML}{B91C1C}
\definecolor{fhegray}{HTML}{4B5563}

\titleformat{\section}{\normalfont\sffamily\large\bfseries\color{fheblue}}{\thesection}{0.7em}{}
\titleformat{\subsection}{\normalfont\sffamily\normalsize\bfseries\color{fheblue!85!black}}{\thesubsection}{0.7em}{}
\titlespacing*{\section}{0pt}{12pt}{5pt}
\titlespacing*{\subsection}{0pt}{8pt}{3pt}

\pagestyle{fancy}\fancyhf{}
\fancyhead[L]{\small\sffamily\color{fhegray}FHE handout}
\fancyhead[R]{\small\sffamily\color{fhegray}Why noise, and what to do about it}
\fancyfoot[C]{\small\sffamily\thepage}
\renewcommand{\headrulewidth}{0.3pt}

%% ---- boxes ------------------------------------------------------------
\newtcolorbox{formal}[1]{enhanced,colback=fheblue!4,colframe=fheblue,boxrule=0.6pt,
  arc=2pt,left=6pt,right=6pt,top=3pt,bottom=3pt,before skip=7pt,after skip=7pt,
  fonttitle=\sffamily\bfseries\small,coltitle=white,colbacktitle=fheblue,title={#1}}
\newtcolorbox{aside}[1][]{enhanced,colback=fheteal!6,colframe=fheteal!6,
  borderline west={2.5pt}{0pt}{fheteal},sharp corners,left=7pt,right=6pt,top=3pt,
  bottom=3pt,before skip=7pt,after skip=7pt,fontupper=\small}
\newtcolorbox{careful}{enhanced,colback=fhered!5,colframe=fhered!5,
  borderline west={2.5pt}{0pt}{fhered},sharp corners,left=7pt,right=6pt,top=3pt,
  bottom=3pt,before skip=7pt,after skip=7pt,fontupper=\small}

%% ---- notation ---------------------------------------------------------
\newcommand{\ZZ}{\mathbb{Z}}
\newcommand{\Zq}{\ZZ_q}
\newcommand{\Enc}{\mathsf{E}}
\newcommand{\Dec}{\mathsf{D}}
\newcommand{\Eval}{\mathsf{V}}
\newcommand{\KeyGen}{\mathsf{K}}
\newcommand{\sk}{\mathit{sk}}
\newcommand{\ek}{\mathit{ek}}
\newcommand{\evk}{\mathit{evk}}
\newcommand{\Msg}{\mathcal{M}}
\newcommand{\Ctx}{\mathcal{C}}
\newcommand{\Fam}{\mathcal{F}}
\newcommand{\hadd}{\boxplus}
\newcommand{\hmul}{\boxtimes}
\renewcommand{\vec}[1]{\bm{#1}}
\newcommand{\inner}[2]{\langle #1,#2\rangle}
\newcommand{\round}[1]{\lfloor #1\rceil}
\newcommand{\indcpa}{\textsf{IND-CPA}}

\begin{document}

%% ---- title -------------------------------------------------------------
\begin{center}
  {\sffamily\bfseries\color{fheblue}\LARGE Computing on Locked Boxes\par}
  \vspace{2pt}
  {\sffamily\large\color{fhegray} A short handout on fully homomorphic encryption\par}
  \vspace{4pt}
  {\color{fheblue}\rule{0.5\textwidth}{0.8pt}}
\end{center}

% \noindent
% Imagine handing a jeweller a locked glovebox full of gold. They can reach in
% through the gloves, cut and solder and polish, but they can never take
% anything out, and they never see the key. At the end you get your box back,
% open it, and find a finished necklace. Fully homomorphic encryption (FHE) is
% that glovebox for data. This handout tells the story of how we got one: what
% the idea is, why it is useful, why the obvious way to build it fails, and why
% every working construction ends up living with something called \emph{noise}.
% It follows Gentry's survey~\cite{gentry} and Halevi's tutorial~\cite{halevi}
% closely.

%% =====================================================================
\section{What is homomorphic encryption?}

Informally, a homomorphic scheme is an ordinary encryption scheme with one
extra public algorithm that computes on ciphertexts. Mathematically it is a
commutative square: operating on ciphertexts and then decrypting gives the
same answer as decrypting and then operating on messages,
\[
\begin{tikzcd}[column sep=large,row sep=normal]
\Ctx\times\Ctx \arrow[r,"\hadd\text{ or }\hmul"] \arrow[d,dashed,"\Dec"'] & \Ctx \arrow[d,dashed,"\Dec"]\\
\Msg\times\Msg \arrow[r,"+\text{ or }\times"] & \Msg
\end{tikzcd}
\]
with the dashed arrows easy only for the key holder. In Gentry's phrase, it is
\emph{a homomorphism that has been hardened}. Now let us say it properly. We
follow Halevi~\cite[\S2.1]{halevi}; $\lambda$ is the security parameter and
$\tau$ a \emph{functionality parameter} that will later bound circuit depth.

\begin{formal}{Definition 1 (Homomorphic encryption: syntax)}
A homomorphic encryption scheme over a plaintext space $\Msg$ is a tuple
$\mathcal{E}=(\KeyGen,\Enc,\Dec,\Eval)$ of probabilistic polynomial-time
algorithms:
\begin{itemize}[leftmargin=*]
\item $(\sk,\ek,\evk)\gets\KeyGen(1^\lambda,1^\tau)$ outputs a secret key, an
  encryption key and a public evaluation key;
\item $c\gets\Enc(\ek,m)$ encrypts $m\in\Msg$ using fresh randomness;
\item $m:=\Dec(\sk,c)$ is deterministic;
\item $c'\gets\Eval(\evk,\Pi,c_1,\dots,c_t)$ takes a circuit $\Pi$ over $\Msg$
  with $t$ inputs and one output, and runs in time $\mathrm{poly}(\lambda,|\Pi|)$.
\end{itemize}
Outputs of $\Enc$ are \emph{fresh} ciphertexts; outputs of $\Eval$ are
\emph{evaluated} ciphertexts. The scheme is \emph{public-key} if $\ek$ is
published and \emph{secret-key} if $\ek=\sk$.
\end{formal}

\begin{formal}{Definition 2 (Correctness for a circuit family)}
Let $\mathcal{C}=\{\mathcal{C}_\tau\}_{\tau\in\mathbb{N}}$ be a family of
circuits. $\mathcal{E}$ is \emph{correct for $\mathcal{C}$}, or
\emph{$\mathcal{C}$-homomorphic}, if for all $\lambda,\tau$ and every key
tuple in the support of $\KeyGen(1^\lambda,1^\tau)$:
\begin{enumerate}[(i),leftmargin=*]
\item $\Dec(\sk,\Enc(\ek,m))=m$ for every $m\in\Msg$, with probability $1$;
\item for every $\Pi\in\mathcal{C}_\tau$ with $t$ inputs and all
  $m_1,\dots,m_t\in\Msg$, with probability $1$,
  \[
  \Dec\bigl(\sk,\ \Eval(\evk,\Pi,c_1,\dots,c_t)\bigr)=\Pi(m_1,\dots,m_t)
  \quad\text{where } c_i\gets\Enc(\ek,m_i).
  \]
\end{enumerate}
\end{formal}

% \begin{formal}{Definition 3 (\indcpa{} security)}
% In the \indcpa{} game the adversary $\mathcal{A}(1^\lambda)$ first outputs
% $1^\tau$; the challenger runs $(\sk,\ek,\evk)\gets\KeyGen(1^\lambda,1^\tau)$
% and gives $\mathcal{A}$ the public keys ($\ek$ if public-key, and $\evk$).
% $\mathcal{A}$ outputs $m_0,m_1\in\Msg$, receives $c^\ast\gets\Enc(\ek,m_b)$ for
% a uniform secret bit $b$, and outputs $b'$. In the secret-key case
% $\mathcal{A}$ also gets an encryption oracle $\Enc(\sk,\cdot)$ throughout.
% $\mathcal{E}$ is \emph{\indcpa{} secure} if for every PPT $\mathcal{A}$,
% $\bigl|\Pr[b'=b]-\tfrac12\bigr|\le\mathrm{negl}(\lambda)$.
% \end{formal}

\noindent
Notice that $\Eval$ does not appear in the security game: it uses no secret, so
the adversary can run it on its own. Two consequences will matter later. The
encryption must be randomized, since otherwise the adversary encrypts $m_0$
itself and compares. And the adversary may collect as many
(ciphertext, message) pairs as it likes. That second fact comes back to bite
us in Section~3.

\paragraph{Examples.} Homomorphic schemes are older than you might think.
\begin{itemize}[leftmargin=*]
\item \textbf{Textbook RSA} ($c=m^e \bmod N$) is multiplicative:
  $m_1^e\,m_2^e=(m_1m_2)^e$. This observation is what made Rivest, Adleman
  and Dertouzos invent the whole idea in 1978. (It is deterministic, though,
  so it is not \indcpa{} secure as is.)
\item \textbf{ElGamal}. In a cyclic group $G=\langle g\rangle$ of prime order
  $q$, the secret is $s$, the public key $h=g^s$, and
  \[
  \Enc(m)=(g^t,\ m\,h^t),\qquad \Dec(c_0,c_1)=c_1\,c_0^{-s}.
  \]
  Multiplying two ciphertexts componentwise gives
  $(g^{t_1+t_2},\ m_1m_2\,h^{t_1+t_2})$, a fresh-looking encryption of
  $m_1m_2$. Security rests on the decisional Diffie--Hellman assumption.
\item \textbf{Paillier} is additive: multiplying ciphertexts in
  $\ZZ_{N^2}^\ast$ adds messages in $\ZZ_N$.
\item \textbf{Goldwasser--Micali} encrypts a bit as a square or non-square
  modulo $N=pq$; multiplying ciphertexts XORs the bits.
\end{itemize}
In the language of Definition~2, each of these is correct for the family of
circuits built from \emph{one} kind of gate, of any size (and $\tau$ plays no
role). A scheme is \emph{fully} homomorphic when the family contains every
Boolean circuit, which needs both XOR and AND. Getting that second gate took
thirty years.

%% =====================================================================
\section{Why would anyone want this?}

The short answer: it lets you outsource computation without outsourcing
trust. A few places where that matters:
\begin{itemize}[leftmargin=*]
\item \textbf{Private cloud computing.} Gentry's own example: you upload
  your encrypted financial records, a tax service fills in your (encrypted)
  return, and you decrypt the result. The service learns nothing.
\item \textbf{Private machine learning.} A hospital can send encrypted
  patient data to a model owner and receive an encrypted diagnosis. Schemes
  designed for approximate arithmetic (CKKS, the ``fourth generation'') are
  aimed precisely at logistic regression and neural networks.
\item \textbf{Encrypted search and private information retrieval.} A client
  can query a database without the server learning which record was asked
  for.
\item \textbf{Secure computation with few messages.} Halevi calls FHE the
  ``Swiss Army knife of cryptography'': two parties can compute a joint
  function with just one message in each direction (given an extra
  property called circuit privacy).
\item \textbf{Sensitive analytics}, such as genomic studies or financial
  risk models across institutions that cannot share raw data.
\end{itemize}
Two honest caveats. FHE is still expensive: Gentry estimates the overhead of
the best current systems as low as about $10{,}000\times$ for suitable
workloads, down from hopeless in 2009. And FHE computations are
\emph{input-oblivious}: the circuit cannot branch on secret data, because
branching would leak it.

%% =====================================================================
\section{Why noise cannot (as far as we know) be avoided}

So we want a ring homomorphism, one that respects both $+$ and $\times$, and
we want to hide it. Let's try the obvious things and watch them break.

\subsection{Attempt one: a linear map, out in the open}

Take a secret vector $\vec s\in\Zq^n$ and let decryption be the inner
product,
\[
\Dec(\vec c)=\inner{\vec c}{\vec s}\bmod q .
\]
It is beautifully homomorphic: adding ciphertexts adds messages, exactly. It
is also completely insecure. Remember the adversary can collect pairs
$(\vec c_i,m_i)$. Each pair is one linear equation in the unknown key:
\[
\inner{\vec c_1}{\vec s}=m_1,\quad \inner{\vec c_2}{\vec s}=m_2,\quad\dots,\quad \inner{\vec c_n}{\vec s}=m_n .
\]
After about $n$ of them, Gaussian elimination hands over $\vec s$. Gentry
remarks in a footnote that a surprising number of published FHE proposals
die for exactly this reason: their decryption is linear.

\begin{aside}
Making decryption merely \emph{non-linear} does not rescue us. If decryption
is a degree-$d$ polynomial in the ciphertext coordinates, treat each monomial
as a new unknown and the same elimination works with about $n^d$ samples.
The real requirement is that decryption be \emph{unlearnable} from examples,
and a rigid ring homomorphism is very easy to learn.
\end{aside}

\subsection{Attempt two: hide the linear map in an exponent}

Here is the clever fix, and it is exactly ElGamal. Keep the inner product, but
compute it in the exponent of a group where discrete logarithms are hard.
Writing $c_0=g^a$, $c_1=g^b$, $m=g^\mu$, ElGamal decryption is
\[
\mu=b-s\,a,
\]
a perfectly linear equation. The attacker just cannot \emph{see} it: to
write down $a,b,\mu$ it would have to take discrete logs. So the linear
algebra attack never gets started, and ElGamal is secure classically.

Wonderful, but now try to multiply. A homomorphic multiplication in the
exponent would be a public map $(g^x,g^y)\mapsto g^{xy}$. Apply it to the
public values and you have solved the computational Diffie--Hellman problem,
which is the very thing ElGamal's security rests on. So
\begin{center}
\emph{the exponent lets you add for free, but multiplying in it would break it.}
\end{center}
(Bilinear pairings are the one partial exception: they buy exactly one
multiplication, landing in a different group. That gives degree-2 schemes
such as Boneh--Goh--Nissim, still far from arbitrary circuits.)

Let's pause on the pattern. Hiding the structure buys security, but only for
\emph{one} operation. FHE needs two.

\subsection{The quantum wall}

There is a second, independent problem, and it does not care how cleverly the
structure is hidden. In any clean scheme, the encryptions of the identity form
a subgroup $H$ of the ciphertext group $G$: the kernel of decryption.
\indcpa{} security demands that random elements of $H$ look like random
elements of $G$.

\begin{formal}{Theorem 4 (Watrous; Armknecht et al.)}
A quantum computer can compute the order of any finite \emph{solvable} group
given by generators. Consequently, a scheme whose ciphertexts form a finite
solvable group, and whose decryption is a group homomorphism, is not
\indcpa{} secure against quantum adversaries.
\end{formal}

\noindent
``Solvable'' means built out of abelian layers: repeatedly taking the
commutator subgroup $G\supseteq [G,G]\supseteq\cdots$ eventually reaches the
identity. Every abelian group qualifies. The attack is simple once you have
order-finding: collect encryptions of the identity, let $K$ be the subgroup
they generate, and test whether the challenge $c$ lies in $K$ by checking
whether adding it changes the order, $|\langle K,c\rangle|=|K|$.

% Now the punchline for FHE. The additive group of \emph{any} ring is abelian,
% and encryptions of $0$ form an ideal, which is a subgroup. So
% \begin{center}
% \emph{every noise-free ring-homomorphic scheme is quantum-insecure, however it is hidden.}
% \end{center}
% GM, ElGamal and Paillier fall to this too (and to Shor), but only against a
% large quantum computer; classically they remain secure.

\subsection{Back to the main road: add noise}

Return to Attempt one and change a single thing. Put a small random error
into every equation:
\[
\inner{\vec c}{\vec s}=m\cdot\round{q/2}+e,\qquad |e|\ll q .
\]
Decryption computes the inner product and rounds to the nearest of $0$ or
$\round{q/2}$. This is Regev's encryption scheme, based on \emph{learning
with errors} (LWE). Watch what happens to the two attacks.
\begin{itemize}[leftmargin=*]
\item \textbf{Gaussian elimination} now solves $A\vec s=\vec b-\vec e$ with
  $\vec e$ unknown. It outputs $\vec s+A^{-1}\vec e$, and for random $A$ the
  term $A^{-1}\vec e$ is uniformly random. Elimination multiplies small errors
  by huge coefficients, and you get a random wrong answer.
\item \textbf{The subgroup} disappears. The set of valid encryptions of $0$ is
  a fuzzy band around the hyperplane $\inner{\vec c}{\vec s}=0$. Adding two of
  them doubles the noise, so the band is not closed, and in fact it generates
  the entire ciphertext space. There is no kernel for Watrous to find.
\end{itemize}
And noise does not care how many operations the underlying homomorphism
respects, so unlike the exponent it hardens $+$ and $\times$ together.

\begin{center}\small
\begin{tabularx}{\textwidth}{@{}l X X l@{}}
\toprule
\textbf{Hardening} & \textbf{Classical} & \textbf{Quantum} & \textbf{Operations}\\
\midrule
none (visible linear) & broken by Gaussian elimination & broken & $+$\\
exponent (ElGamal, \dots) & secure under DDH, QR, DCR & broken (Shor, Watrous) & one\\
noise (LWE) & secure under LWE & no known attack & $+$ and $\times$\\
\bottomrule
\end{tabularx}
\end{center}

\begin{careful}
\textbf{A word of honesty.} ``Cannot be avoided'' should really be ``nobody
knows how to avoid it''. Noise-free FHE is an open problem (Gentry \S5).
Quantumly, noise-free ring schemes are ruled out by the theorem above, and
the only clean escape left is a \emph{non-solvable} group such as $A_5$, for
which no secure candidate is known. Classically, there is no impossibility
theorem at all; it is simply that every concrete attempt has been broken.
\end{careful}

%% =====================================================================
\section{A toy example: watching the noise blow up}

Noise buys security, but it has a price, and the price is easiest to see in
the integer scheme of van Dijk, Gentry, Halevi and Vaikuntanathan~\cite{dghv}.
It rests on the \emph{approximate GCD} problem: given many numbers that are
close to multiples of a secret odd $p$, find $p$.

\begin{formal}{A symmetric toy scheme over the integers}
Secret key: an odd integer $p$. To encrypt a bit $m$, pick a random $q$ and a
small random $r$ and output
\[
c=q\,p+\underbrace{2r+m}_{\text{noise term}} .
\]
To decrypt, reduce $c$ modulo $p$ into $(-p/2,p/2]$ and take the result
modulo $2$. This works as long as the noise term stays below $p/2$ in
absolute value.
\end{formal}

% \noindent
% Take $p=1009$. Encrypt $m_1=1$ with $q=17$, $r=3$, and $m_2=0$ with $q=23$,
% $r=5$:
% \[
% c_1=17\cdot1009+7=17160,\qquad c_2=23\cdot1009+10=23217 .
% \]
% \textbf{Addition.} $c_1+c_2=40377$, and $40377\bmod 1009=17$, which is odd,
% so it decrypts to $1=m_1\oplus m_2$. The noise went from $7$ and $10$ to
% $17$: noises \emph{add}.

% \textbf{Multiplication.} $c_1c_2=398403720$, and modulo $1009$ this is
% $70=7\times10$, which is even, giving $0=m_1\wedge m_2$. Correct, but the noise
% went from $7$ and $10$ to $70$: noises \emph{multiply}.

% So far so good. Now take a fresh encryption of $1$ with noise term $9$
% ($q=12$, $r=4$, so $c=12117$) and square it twice, which should still encrypt
% $1\cdot1\cdot1\cdot1=1$:
% \[
% \begin{array}{lll}
% c & \text{noise } 9 & \text{decrypts to } 1 \ \checkmark\\
% c^2 & \text{noise } 81 & \text{decrypts to } 1 \ \checkmark\\
% c^4 & \text{noise } 6561 \equiv -502 \pmod{1009} & \text{decrypts to } 0 \ \boldsymbol{\times}
% \end{array}
% \]
% The true noise $6561$ has passed $p/2\approx504$, so reducing modulo $p$
% wraps it around to $-502$, and the parity flips. The ciphertext is now
% garbage, through no fault of the computation.

\begin{aside}
In general, if fresh noise is about $B$, then after $d$ levels of
multiplication it is about $B^{2^d}$, while additions only make it grow
linearly. With LWE the picture is the same in spirit: Gentry shows each
multiplication multiplies the noise by a polynomial factor $\mathrm{poly}(n)$,
so after depth $d$ the noise is about $\varepsilon_0\cdot\mathrm{poly}(n)^d$.
Either way, \textbf{multiplicative depth is the budget}, and it runs out.
\end{aside}

%% =====================================================================
\section{The way forward: from levels to full FHE}

The toy example leaves us with a scheme that is correct up to some depth and
no further. So the natural thing is to make the depth part of the contract.
Before we do, one more requirement, without which every definition below
would be empty.

\begin{formal}{Definition 5 (Compactness, Halevi Def.~7)}
$\mathcal{E}$ is \emph{compact} if there is a fixed polynomial $B(\cdot)$ such
that for all $\lambda,\tau$, every single-output circuit $\Pi$ and all inputs,
every evaluated ciphertext $c'\gets\Eval(\evk,\Pi,c_1,\dots,c_t)$ satisfies
$|c'|\le B(\lambda)$, independently of $\tau$ and of $\Pi$.
\end{formal}

\noindent
Why insist on this? Consider the ``scheme'' whose $\Eval$ simply outputs the
pair $(\Pi,(c_1,\dots,c_t))$ and whose decryption decrypts every $c_i$ and then
runs $\Pi$ itself. It is correct for \emph{every} circuit and as secure as the
underlying encryption, and it is useless: the key holder does all the work.
Compactness rules it out by forcing the result of a computation to be as
short as a fresh ciphertext.

Now the two definitions this handout has been heading towards. Fix a
functionally complete gate set, say $\{\mathrm{XOR},\mathrm{AND}\}$ over
$\Msg=\{0,1\}$ (or just $\{\mathrm{NAND}\}$), and write $\mathrm{depth}(\Pi)$
for the length of the longest input-to-output path in $\Pi$.

\begin{formal}{Definition 6 (Leveled fully homomorphic encryption)}
$\mathcal{E}$ is \emph{leveled fully homomorphic} if it is compact and correct
for the family $\mathcal{C}=\{\mathcal{C}_\tau\}$ where
\[
\mathcal{C}_\tau=\{\,\Pi \text{ Boolean circuit}:\ \mathrm{depth}(\Pi)\le\tau\,\}.
\]
The parameter $\tau=L$ is the \emph{number of levels}. The keys output by
$\KeyGen(1^\lambda,1^L)$ may grow with $L$ (polynomially), but evaluated
ciphertexts may not: their length is bounded by $B(\lambda)$ alone.
\end{formal}

\noindent
Gentry's version (Def.~12 of~\cite{gentry}) is the same idea phrased with a
key for every level: $\KeyGen$ produces $(\sk^{(i)},\ek^{(i)},\evk^{(i)})$ for
$i=0,\dots,L$, each gate moves a ciphertext from level $i$ to level $i-1$, and
the algorithms are otherwise independent of $L$. Many practical schemes in the
BGV family~\cite{bgv} meet a relaxed variant, \emph{weak compactness}, where
the ciphertext length may grow polynomially in $L$ but must stay smaller
than the circuits being evaluated (Halevi Def.~8).

\begin{formal}{Definition 7 (Fully homomorphic encryption)}
$\mathcal{E}$ is \emph{fully homomorphic} if it is compact and correct for a
family with
\[
\mathcal{C}_1\ \supseteq\ \{\text{all Boolean circuits, of every depth and size}\}.
\]
In this case $\tau$ is irrelevant and we write $\KeyGen(1^\lambda)$: a single
key tuple serves every computation.
\end{formal}

\noindent
Let's pause and look at the difference, because it is entirely a matter of
quantifier order:
\[
\begin{array}{lll}
\textbf{Leveled FHE:} & \forall L\ \ \exists\,(\sk,\ek,\evk)_L & \forall\,\Pi \text{ with } \mathrm{depth}(\Pi)\le L:\ \text{correct},\\[2pt]
\textbf{FHE:}         & \exists\,(\sk,\ek,\evk)\ \ \forall L & \forall\,\Pi \text{ with } \mathrm{depth}(\Pi)\le L:\ \text{correct}.
\end{array}
\]
A leveled scheme asks you to name the depth before you see a key; an FHE
scheme lets the levels be unbounded. In particular every FHE scheme is
trivially leveled (ignore $L$), while turning a leveled scheme into a fully
homomorphic one is the genuinely hard direction.

\begin{aside}
Gentry phrases FHE as an \emph{unbounded} commutative diagram (Def.~11
of~\cite{gentry}): there is a set of ciphertexts $\Ctx$, containing the fresh
ones, that is \emph{closed} under the homomorphic gates and on which the
square of Section~1 commutes. Evaluated ciphertexts can then be fed back into
$\Eval$ forever, and their size and decryption cost never depend on how many
gates produced them. In a leveled scheme the corresponding set $\Ctx^{(i)}$ is
only mapped to $\Ctx^{(i-1)}$, and after level $0$ the road ends.
\end{aside}

\begin{center}\small
\begin{tabularx}{\textwidth}{@{}l X X@{}}
\toprule
\textbf{Notion} & \textbf{Correct for} & \textbf{Example} \\
\midrule
single-gate, unbounded & circuits of one gate type, any size & GM (XOR), ElGamal ($\times$), Paillier ($+$) \\
somewhat / bounded & circuits of small, fixed depth & the toy integer scheme of Section~4 \\
leveled FHE & all circuits of depth $\le L$, $L$ fixed at key generation & LWE-based BGV, GSW with $q$ chosen for $L$ \\
FHE & all circuits, one key & needs a further idea \\
\bottomrule
\end{tabularx}
\end{center}

\paragraph{Where this leaves us.} The story has a slightly ironic shape. We
wanted an exact ring homomorphism, and exactness is precisely what got us
broken. So we blurred it with noise, accepted that the noise grows, and
discovered that the honest thing to promise is a \emph{level budget}: leveled
FHE. Removing that budget, so that one key carries us through unboundedly many
levels, is what turns leveled FHE into full FHE, and it is the next chapter of
the story.

\begin{thebibliography}{9}\small
\bibitem{gentry} C.~Gentry, \emph{Homomorphic encryption: a mathematical
  survey}, Proc.\ ICM 2022, vol.~2, 956--1006.
\bibitem{halevi} S.~Halevi, \emph{Homomorphic encryption}, in Tutorials on
  the Foundations of Cryptography, Springer, 2017, 219--276.
\bibitem{dghv} M.~van Dijk, C.~Gentry, S.~Halevi, V.~Vaikuntanathan,
  \emph{Fully homomorphic encryption over the integers}, EUROCRYPT 2010.
\bibitem{bgv} Z.~Brakerski, C.~Gentry, V.~Vaikuntanathan, \emph{(Leveled)
  fully homomorphic encryption without bootstrapping}, ACM TOCT 6(3), 2014.
\end{thebibliography}

\end{document}